%%%% RESUMO
%%
%% Apresentação concisa dos pontos relevantes de um texto, fornecendo uma visão rápida e clara do conteúdo e das conclusões do
%% trabalho.

\begin{resumoutfpr}%% Ambiente resumoutfpr

In software development, metrics are commonly used to assess code quality and development time. Professionally, this is often done via pull requests, whereas in academia, assessment relies on occasional submissions through emails or Learning Management Systems (LMS). Although LMSs record activity logs, a significant gap remains in the structured collection of data on the programming learning process, particularly in computer lab environments. This work proposes creating a dataset of code development logs to capture the learning trail in real-time. Unlike conventional approaches that rely on LMSs, data will be collected via a plugin integrated into a widely used code editor. The resulting database is expected to help identify student difficulties, enable personalized teaching, and improve pedagogical practices.

\end{resumoutfpr}

%De acordo com a NBR 6028:2021, a apresentação gráfica deve seguir o padrão do documento no qual o resumo está inserido. Para definição das palavras-chave (e suas correspondentes em inglês no abstract) consultar em Termo tópico do Catálogo de Autoridades da Biblioteca Nacional, disponível em: http://acervo.bn.gov.br/sophia_web/autoridade